% !TeX root = ../../Book3.tex
% 纲要标题：### 8.4 可验证的计算与系数域
% 纲要位置：工作记录/计划纲要.md，第 2061 行。
% 纲要细目（写作范围）：
% * 有理数、有限域与有理函数域上的精确运算
% * 输入、输出、正确性证书及可重复计算的记录格式
% * 终止性与实际复杂度的区别；系数膨胀和单项式序的影响
% * 数值近似不能替代理想成员与理想相等的精确证明


\section{可验证的计算与系数域}
\label{b3:sec:0804}


同一组符号式放在不同的系数环中，可能给出不同的理想。精确计算除了记录多项式，还须说明允许哪些除法、参数取值是否使分母为零，以及怎样独立检验结果。本节用具体算例说明这些要求，并区分算法终止、计算规模和数值近似。

\subsection{有理数、有限域与有理函数域上的精确运算}
\label{b3:subsec:0804:01}

前面各算法只用到系数的加、减、乘、除及零判定。因此，给定一个抽象域还不够：实际执行时，必须说明它的元素怎样表示、这些运算怎样完成。\(\mathbb Q\) 中的系数可用互素整数对表示，每次运算后约分；\(\mathbb F_p\) 中的系数用模 \(p\) 的剩余类表示，除法由扩展 Euclid 算法求逆。对于有限扩域 \(\mathbb F_p[z]/(h)\)，还须给定不可约多项式 \(h\)，并把系数约化为次数小于 \(\deg h\) 的多项式。

\ProseParagraphBreak
在 \(k(a_1,\ldots,a_s)\) 上计算时，参数 \(a_i\) 被视为代数无关元。系数是多项式之比，零判定归结为分子是否为零多项式；其分母在有理函数域中可逆，却未必在参数的每个具体取值处可逆。例如 \((ax-1)\subset k(a)[x]\) 的约化 Gröbner 基为 \(\{x-a^{-1}\}\)。当 \(a=c\ne0\) 时它可以直接专门化；当 \(a=0\) 时原生成元成为 \(-1\)，所得理想是单位理想，原公式则没有意义。

\begin{proposition}\label{b3:prop:groebner-specialization-open}
设 \(k\) 为域，\(A=k[a_1,\ldots,a_s]\)、\(K=\operatorname{Frac}(A)\)，并固定 \(K[x_1,\ldots,x_n]\) 中的全局单项式序。若 \(F\subset A[x_1,\ldots,x_n]\) 为有限集，\(G\) 是 \((F)K[x_1,\ldots,x_n]\) 的首一 Gröbner 基，则存在非零 \(d\in A\)，使 \(G\subset A_d[x_1,\ldots,x_n]\)，且对任意域扩张 \(k\subset L\) 和任意满足 \(d(c)\ne0\) 的 \(c\in L^s\)，专门化后的 \(G(c)\) 是 \(\Bigl(F(c)\Bigr)\) 的 Gröbner 基，其首单项式与 \(G\) 相同。
\end{proposition}
\begin{proof}
记录 \(G\) 的全部系数、两组生成元互相表示的系数，以及每个 \(S(g_i,g_j)\) 约化为零所得的有界表示系数。这些都是有限多个有理函数。取其所有非零分母的乘积为 \(d\)，便可把全部恒等式写在 \(A_d[x_1,\ldots,x_n]\) 中。

当 \(d(c)\ne0\) 时，代入给出良定义的环同态 \(A_d\to L\)。两组互相表示的等式保证 \(\Bigl(G(c)\Bigr)=\Bigl(F(c)\Bigr)\)。每个 \(g_i\) 首一，故专门化不改变其首单项式；其余系数即使变为零，也只会删去较低的项。因此，各 \(S\)-多项式专门化后仍具有小于相应最小公倍数的有界表示。由\cref{b3:thm:buchberger-criterion}，\(G(c)\) 是所求 Gröbner 基。
\end{proof}

这一命题只说明存在一个可以沿用计算结果的非空开集。对 \(d(c)=0\) 的参数必须另算，不能把它们当作无解。例如 \(F=\{ax,y-x\}\) 在 \(k(a)[x,y]\) 中生成 \((x,y)\)，而在 \(a=0\) 时生成 \((y-x)\)，零点从一点变成一条直线。按参数分情形的系统方法留待附录B。

\subsection{计算步骤与结果验证}
\label{b3:subsec:0804:02}

一份可重复的计算至少应写出系数域、变量及其顺序、单项式序、原生成元，以及输出多项式的精确系数。若使用有限扩域，应写出定义该域的不可约多项式；若使用有理函数域，应区分参数和多项式变量。例如，\(k(a)[x]\) 中可用 \(a\) 除，而 \(k[a,x]\) 中一般不可以，这会改变理想成员问题本身。

\ProseParagraphBreak
验证 Gröbner 基时可把任务分成三个有限检查。给原生成元 \(F=(f_i)\) 和候选基 \(G=(g_j)\)，先检查多项式恒等式
\[
g_j=\sum_i a_{ij}f_i,\qquad f_i=\sum_j b_{ji}g_j,
\]
证明两者生成同一理想；再检查每一对 \(g_i,g_j\) 的 \(S\)-多项式有\cref{b3:thm:buchberger-criterion}要求的有界表示，或按已指定的除法次序确实得到零余式。前两项不能由第三项取代：\(\{x\}\) 自身是 Gröbner 基，却不是 \((x^2)\) 的 Gröbner 基。

\ProseParagraphBreak
理想成员的证书可以更短。证明 \(f\in(F)\)，只需给出并核验 \(f=\sum a_if_i\)，无须重新验证一次完整 Gröbner 基计算；证明 \(f\notin(F)\)，则可以先验证 \(G\) 的 Gröbner 基资格，再展示 \(f\) 的非零正规形。保存计算过程中的表示向量，正是为了在结果出来后提取这些较短的恒等式。它们能由另一套程序或人工逐项检查，与发现结果所用的实现相互独立。

\subsection{复杂度、系数膨胀与单项式序}
\label{b3:subsec:0804:03}

Buchberger 算法的终止性来自首单项式理想不可能形成无限严格升链。这个论证没有给出“只需少量步骤”的保证，也没有限制途中多项式的次数、项数和系数长度。即使最后得到很短的答案，中间表达式也可能很大。

\ProseParagraphBreak
消去本身便可能提高次数。在 \(k[x_1,\ldots,x_n]\) 中考虑
\[
I=(x_2-x_1^2,x_3-x_2^2,\ldots,x_n-x_{n-1}^2).
\]
逐次代入表明商环同构于 \(k[x_1]\)，且
\[
I\cap k[x_1,x_n]=(x_n-x_1^{2^{n-1}}).
\]
后一等式也可直接检验：以 \(x_n-x_1^{2^{n-1}}\) 作一元首一除法，余式只含 \(x_1\)，代入后为零便迫使余式为零。原方程全是二次式，消去后的关系却有指数增长的次数。这是具体输出变大的例子，并不声称所有单项式序或所有表示方法都具有相同复杂度。

\ProseParagraphBreak
系数也可能增长。对方程 \(x_1=2\)、\(x_{i+1}=x_i^2\) 逐次消去，得到整数 \(x_n=2^{2^{n-1}}\)，其二进制长度为 \(2^{n-1}+1\)。有理数运算虽然精确，存储分子、分母和做整数运算的成本仍须计算。单项式序还会改变中间基及最后基的形状：字典序适合消去，但不总是最快；先用另一种全局序得到商代数，再在线性代数中换序，是下一节要解释的做法。附录B再讨论具体实现，正文只以已经证明终止的算法承担存在性与有效性。

\subsection{数值近似与精确理想判定}
\label{b3:subsec:0804:04}

在 \(\mathbb R[x]\) 中，\(I_\varepsilon=(x,x-\varepsilon)\) 当 \(\varepsilon\ne0\) 时等于整个环，因为
\[
1=\varepsilon^{-1}x-\varepsilon^{-1}(x-\varepsilon).
\]
当 \(\varepsilon=0\) 时却等于 \((x)\)。若把足够小的非零系数统一舍为零，便会把一个无公共零点的系统误判为有零点的系统。这里出错的不是近似精度还不够高，而是“系数等于零”这个精确问题没有被近似数据决定。

\ProseParagraphBreak
同样，在 \(\mathbb C[x]\) 中 \(x^2\) 与 \(x^2-\varepsilon\) 的系数可以任意接近，但前者有一个二重根，后者在 \(\varepsilon\ne0\) 时有两个不同的根。数值计算适合近似根的位置，精确理想运算则要判断恒等式、理想相等和重数。两者可以相互帮助：例如先猜测有理系数关系，再用精确运算验证；但未经验证的近似零余式不能作为理想成员关系的证明。
